\newpage
\section{Panjang Vektor}

\begin{outcome}

\begin{enumerate}

\item[A.] Menentukan panjang vektor dan jarak antara dua titik dalam $\mathbb{R}^n$.

\item[B.] Menentukan vektor satuan yang bersesuaian dengan suatu vektor dalam $\mathbb{R}^n$.

\end{enumerate}
\end{outcome}

Dalam bagian ini, kita mengkaji apa yang dimaksud dengan panjang suatu vektor dalam $\mathbb{R}^n$.
Kita mengembangkan konsep ini dengan terlebih dahulu meninjau jarak antara dua titik dalam $\mathbb{R}^n$.

Pertama, kita akan membahas konsep jarak untuk $\mathbb{R}$, yaitu untuk titik-titik dalam $\mathbb{R}^1$.
Di sini,
jarak antara dua titik $P$ dan $Q$ diberikan oleh nilai mutlak
selisih keduanya. Kita menyatakan jarak antara $P$ dan $Q$ dengan $d(P,Q)$, yang didefinisikan sebagai
\begin{equation}
d(P,Q) = \sqrt{ \left( P-Q\right) ^{2}}
\label{distance1}
\end{equation}

Sekarang perhatikan kasus $n=2$, yang diperagakan oleh gambar berikut.

\begin{picture}(1,100)
\put(100,20){\begin{picture}(1,1)
\setlength{\unitlength}{2pt}
\thicklines
\qbezier[10](0,0)(25,0)(50,0)
\qbezier[7](0,0)(0,15)(0,30)
\qbezier[10](0,30)(25,30)(50,30)
\qbezier[7](50,30)(50,15)(50,0)
\qbezier(0,0)(25,15)(50,30)
\put(0,0){\circle*{2}}
\put(50,0){\circle*{2}}
\put(50,30){\circle*{2}}
\put(5,-5){$Q=(q_1,q_2)$}
\put(55,-5){$(p_1,q_2)$}
\put(55,33){$P=(p_1,p_2)$}
\end{picture}}
\end{picture}\par\captionof{figure}{Jarak antara dua titik pada bidang.}\label{fig:distance-2d}% fig-num-auto


Terdapat dua titik $P =\left( p_{1},p_{2}\right) $ dan
$Q = \left(q_{1},q_{2}\right)$ pada bidang. Jarak antara kedua
titik ini ditunjukkan pada gambar sebagai garis utuh. Perhatikan bahwa garis ini
merupakan hipotenusa suatu segitiga siku-siku yang merupakan separuh dari persegi panjang
yang ditunjukkan dengan garis putus-putus. Kita ingin mencari panjang hipotenusa ini,
yang akan memberikan jarak antara kedua titik tersebut. Perhatikan bahwa
panjang sisi-sisi segitiga ini adalah $\left| p_{1}-q_{1}\right| $
dan $\left| p_{2}-q_{2}\right|$, yaitu nilai mutlak selisih nilai-nilai tersebut. Oleh karena itu, Teorema Pythagoras
menyiratkan bahwa panjang hipotenusa (dan dengan demikian jarak antara $P$ dan $Q$) sama dengan
\begin{equation}
\left( \left| p_{1}-q_{1}\right| ^{2}+\left| p_{2}-q_{2}\right| ^{2}\right)
^{1/2}=\left( \left( p_{1}-q_{1}\right) ^{2}+\left( p_{2}-q_{2}\right)
^{2}\right) ^{1/2}
\label{distance2}
\end{equation}

Sekarang misalkan $n=3$ dan misalkan $P = \left( p_{1},p_{2},p_{3}\right) $ dan $Q = \left(
q_{1},q_{2},q_{3}\right) $ adalah dua titik dalam $\mathbb{R}^{3}.$ Perhatikan
gambar berikut, dengan garis utuh yang menghubungkan kedua titik dan sebuah
garis putus-putus yang menghubungkan titik $\left( q_{1},q_{2},q_{3}\right) $ dan $\left(
p_{1},p_{2},q_{3}\right) .$ 

\begin{picture}(1,139)
\put(100,20){\begin{picture}(1,1)
\setlength{\unitlength}{2pt}
\thicklines
\qbezier[10](0,0)(25,0)(50,0)
\qbezier[7](0,0)(0,15)(0,30)
\qbezier[10](0,30)(25,30)(50,30)
\qbezier[7](50,30)(50,15)(50,0)
\put(20,20){
\qbezier[10](0,0)(25,0)(50,0)
\qbezier[7](0,0)(0,15)(0,30)
\qbezier[10](0,30)(25,30)(50,30)
\qbezier[7](50,30)(50,15)(50,0)}
\qbezier[7](0,0)(10,10)(20,20)
\put(0,0){\circle*{2}}
\put(50,0){\circle*{2}}
\put(0,30){\qbezier[7](0,0)(10,10)(20,20)\put(70,20){\circle*{2}}\put(70,-10){\circle*{2}}}
\put(50,30){\qbezier[7](0,0)(10,10)(20,20)}
\put(50,0){\qbezier[7](0,0)(10,10)(20,20)}
\qbezier[13](0,0)(35,10)(70,20)
\qbezier(0,0)(35,25)(70,50)
\put(-30,-5){$Q=(q_1,q_2,q_3)$}
\put(35,-5){$(p_1,q_2,q_3)$}
\put(72,18){$(p_1,p_2,q_3)$}
\put(70,52){$P=(p_1,p_2,p_3)$}
\end{picture}}
\end{picture}\par\captionof{figure}{Jarak antara dua titik dalam tiga dimensi.}\label{fig:distance-3d}% fig-num-auto


Di sini, kita perlu menggunakan Teorema Pythagoras dua kali untuk menemukan panjang
garis utuh tersebut. Pertama, berdasarkan Teorema Pythagoras, panjang garis putus-putus yang menghubungkan $\left(
q_{1},q_{2},q_{3}\right) $ dan $\left( p_{1},p_{2},q_{3}\right) $ sama dengan
\begin{equation*}
\left( \left( p_{1}-q_{1}\right) ^{2}+\left( p_{2}-q_{2}\right) ^{2}\right)
^{1/2}
\end{equation*}
sedangkan panjang garis yang menghubungkan $\left( p_{1},p_{2},q_{3}\right) $ dengan
$\left( p_{1},p_{2},p_{3}\right) $ hanyalah $\left| p_{3}-q_{3}\right| .$ Oleh karena itu, dengan menerapkan Teorema Pythagoras sekali lagi, panjang garis yang menghubungkan
titik-titik $P = \left( p_{1},p_{2},p_{3}\right) $ dan $Q = \left(
q_{1},q_{2},q_{3}\right) $ sama dengan
\begin{equation*}
\left( \left( \left( \left( p_{1}-q_{1}\right) ^{2}+\left(
p_{2}-q_{2}\right) ^{2}\right) ^{1/2}\right) ^{2}+\left( p_{3}-q_{3}\right)
^{2}\right) ^{1/2}
\end{equation*}
\begin{equation}
 =\left( \left( p_{1}-q_{1}\right) ^{2}+\left( p_{2}-q_{2}\right)
^{2}+\left( p_{3}-q_{3}\right) ^{2}\right) ^{1/2}
\label{distance3}
\end{equation}

Pembahasan ini memotivasi definisi berikut untuk jarak antara titik-titik dalam $\mathbb{R}^n$.

\begin{definition}{Jarak antara Titik-Titik}\label{def:distancebetweenpoints}
Misalkan $P=\left( p_{1},\cdots ,p_{n}\right) $ dan
$Q=\left( q_{1},\cdots ,q_{n}\right) $ adalah dua titik dalam
$\mathbb{R}^{n}$. Maka jarak
antara kedua titik ini didefinisikan sebagai
\begin{equation*}
\text{ jarak antara }P\text{ dan } Q\text{ } =
d( P, Q ) = \left( \sum_{k=1}^{n}\left\vert
p_{k}-q_{k}\right\vert ^{2}\right) ^{1/2}
\end{equation*}
Ini disebut \textbf{rumus jarak}. Kita juga dapat menulis $\left\vert P - Q \right\vert $ sebagai jarak antara $P$ dan $Q$.
\index{rumus jarak}
\end{definition}

Dari pembahasan di atas, Anda dapat melihat bahwa Definisi \ref{def:distancebetweenpoints} berlaku untuk kasus-kasus khusus $n=1,2,3$, seperti dalam
Persamaan \ref{distance1}, \ref{distance2}, \ref{distance3}.
Dalam contoh berikut, kita menggunakan Definisi \ref{def:distancebetweenpoints} untuk menemukan jarak antara dua titik dalam
$\mathbb{R}^4$.

\begin{example}{Jarak antara Titik-Titik}\label{exa:distancebetweenpoints}
Tentukan jarak antara titik $P$ dan $Q$ dalam $\mathbb{R}^{4}$,
dengan $P$ dan $Q$ diberikan oleh
\begin{equation*}
P=\left( 1,2,-4,6\right)
\end{equation*}
dan
\begin{equation*}
Q=\left( 2,3,-1,0\right)
\end{equation*}
\end{example}

\begin{solution}
Kita akan menggunakan rumus yang diberikan dalam Definisi \ref{def:distancebetweenpoints} untuk mencari jarak antara
$P$ dan $Q$.
Gunakan rumus jarak dan tuliskan
\begin{equation*}
d(P,Q)= \left( \left( 1-2\right) ^{2}+\left( 2-3\right)
^{2}+\left( -4-\left( -1\right) \right) ^{2}+\left( 6-0\right)^{2}\right) ^{\frac{1}{2}} = 47
\end{equation*}

Oleh karena itu, $d( P,Q) =
\sqrt{47}.$

\end{solution}

Terdapat sifat-sifat tertentu dari jarak antara titik-titik yang penting dalam kajian kita.
Sifat-sifat ini diuraikan dalam teorema berikut.

\begin{theorem}{Sifat-Sifat Jarak}\label{thm:distanceproperties}
Misalkan $P$ dan $Q$ adalah titik-titik dalam $\mathbb{R}^n$, dan misalkan jarak di antaranya,
$d( P, Q)$, diberikan seperti dalam Definisi \ref{def:distancebetweenpoints}.
Maka sifat-sifat berikut berlaku \index{rumus jarak!sifat-sifat}.
\begin{itemize}
\item $ d( P, Q) = d( Q, P)  $
\item $ d( P, Q) \geq 0$, dan sama dengan 0 tepat ketika $P = Q.$
\end{itemize}
\end{theorem}

Konsep jarak memiliki banyak penerapan. Misalnya,
dengan diberikan dua titik, kita dapat menanyakan kumpulan titik apa yang
semuanya berjarak sama dari kedua titik yang diberikan tersebut. Hal ini dikaji dalam
contoh berikut.

\begin{example}{Bidang di Antara Dua Titik}\label{exa:planebetweentwopoints}
Jelaskan titik-titik dalam $\mathbb{R}^3$ yang berjarak sama dari $\left(
1,2,3\right) $ dan $\left( 0,1,2\right) .$ 
\end{example}

\begin{solution}
Misalkan $P = \left( p_1 , p_2, p_3\right) $ adalah titik semacam itu. Oleh karena itu, $P$ berjarak sama dari $\left(
1,2,3\right) $ dan $\left( 0,1,2\right) .$ 
Kemudian, berdasarkan Definisi \ref{def:distancebetweenpoints},
\begin{equation*}
\sqrt{\left( p_1 -1\right) ^{2}+\left( p_2 -2\right) ^{2}+\left( p_3-3\right) ^{2}}=
\sqrt{\left( p_1 - 0 \right)^{2}+\left( p_2-1\right) ^{2}+\left( p_3-2\right) ^{2}}
\end{equation*}
Dengan menguadratkan kedua ruas, kita memperoleh
\begin{equation*}
\left( p_1 -1\right) ^{2}+\left( p_2 -2\right) ^{2}+\left( p_3 -3\right)
^{2}=p_1^{2}+\left( p_2-1\right) ^{2}+\left( p_3 -2\right) ^{2}
\end{equation*}
sehingga
\begin{equation*}
\allowbreak p_1^{2}-2p_1+14+p_2^{2}-4p_2+p_3^{2}-6p_3=p_1^{2}+p_2^{2}-2p_2+5+p_3^{2}-4p_3
\end{equation*}
Dengan menyederhanakannya, diperoleh
\begin{equation*}
-2p_1+14-4p_2-6p_3=-2p_2+5-4p_3
\end{equation*}
yang dapat ditulis sebagai
\begin{equation}
2p_1+2p_2+2p_3=-9  \label{distanceplane}
\end{equation}
Oleh karena itu, titik-titik $P = \left(
p_1,p_2,p_3\right) $ yang berjarak sama
dari masing-masing titik yang diberikan membentuk suatu bidang yang persamaannya diberikan oleh \ref{distanceplane}.
\end{solution}

Sekarang kita dapat menggunakan pemahaman kita tentang jarak antara dua titik untuk mendefinisikan apa yang dimaksud dengan panjang suatu
vektor. Perhatikan definisi berikut.

\begin{definition}{Panjang Vektor}\label{def:lengthofvector}
Misalkan $\vect{u} = \leftB u_{1} \cdots u_{n} \rightB^T$ adalah sebuah vektor dalam
$\mathbb{R}^n$. Maka panjang $\vect{u}$, yang ditulis $\vectlength \vect{u} \vectlength$ \index{vektor!panjang}, diberikan oleh
\begin{equation*}
\vectlength
\vect{u}
\vectlength
= \sqrt{ u_{1}^2 + \cdots + u_{n}^2}
\end{equation*}
\end{definition}

Definisi ini bersesuaian dengan Definisi
\ref{def:distancebetweenpoints}, jika Anda menganggap vektor $\vect{u}$
berpangkal di titik $0 = \left( 0, \cdots ,0 \right)$ dan
berujung di titik $U = \left(u_1, \cdots, u_n \right)$. Maka panjang
$\vect{u}$ sama dengan jarak antara $0$ dan $U$, yaitu $d(0,U)$. Secara umum, $d(P,Q) = \vectlength \longvect{PQ} \vectlength$.

Perhatikan Contoh \ref{exa:distancebetweenpoints}. Berdasarkan Definisi \ref{def:lengthofvector}, kita juga dapat mencari jarak antara $P$ dan $Q$
sebagai panjang vektor yang menghubungkan keduanya. Jadi, jika kita menggambar vektor $\longvect{PQ}$
dengan pangkal di $P$ dan ujung di $Q$, vektor ini akan memiliki panjang
$\sqrt{47}$.

Kita mengakhiri bagian ini dengan definisi baru untuk kasus khusus
vektor yang panjangnya $1$.

\begin{definition}{Vektor Satuan}\label{def:unitvector}
Misalkan $\vect{u}$ adalah sebuah vektor dalam $\mathbb{R}^{n}$. Maka kita menyebut $\vect{u}$ sebagai
\textbf{vektor satuan} \index{vektor satuan} jika panjangnya 1, yaitu jika
\begin{equation*}
\vectlength \vect{u} \vectlength
=
1
\end{equation*}
\end{definition}

Misalkan $\vect{v}$ adalah sebuah vektor dalam $\mathbb{R}^{n}$. Maka vektor $\vect{u}$
yang searah dengan $\vect{v}$ tetapi panjangnya sama dengan $1$ adalah vektor satuan yang bersesuaian dengan $\vect{v}$.
Vektor ini diberikan oleh
\begin{equation*}
\vect{u}
=
\frac{1}{\vectlength \vect{v} \vectlength}
\vect{v}
\end{equation*}

Kita sering menggunakan istilah \textbf{menormalkan} untuk merujuk pada proses ini. Ketika kita \textbf{menormalkan} sebuah vektor, kita mencari vektor satuan yang bersesuaian dan panjangnya $1$.
Perhatikan contoh berikut.

\begin{example}{Menentukan Vektor Satuan}\label{exa:unitvector}
Misalkan $\vect{v}$ diberikan oleh
\begin{equation*}
\vect{v}
=
\leftB
\begin{array}{rrr}
1 & -3 & 4
\end{array}
\rightB^T
\end{equation*}
Tentukan vektor satuan $\vect{u}$ yang searah dengan $\vect{v}$ \index{vektor!vektor satuan yang bersesuaian}.
\end{example}

\begin{solution}
Kita akan menggunakan Definisi \ref{def:unitvector} untuk menyelesaikan soal ini.
Oleh karena itu, kita perlu mencari panjang $\vect{v}$ yang, berdasarkan Definisi \ref{def:lengthofvector},
diberikan oleh
\begin{equation*}
\vectlength
\vect{v}
\vectlength
= \sqrt{ v_{1}^2 + v_{2}^2+ v_{3}^2}
\end{equation*}
Dengan menggunakan nilai-nilai yang bersesuaian, kita memperoleh
\begin{eqnarray*}
\vectlength
\vect{v}
\vectlength
&=& \sqrt{ 1^2 + \left(-3 \right)^2 + 4^2} \\
&=& \sqrt{ 1 + 9 + 16} \\
&=& \sqrt{26}
\end{eqnarray*}
Untuk mencari $\vect{u}$, kita membagi $\vect{v}$ dengan $\sqrt{26}$.
Hasilnya adalah
\begin{eqnarray*}
\vect{u}
&=&
\frac{1}{\vectlength \vect{v} \vectlength}
\vect{v} \\
&=&
\frac{1}{\sqrt{26}}
\leftB
\begin{array}{rrr}
1 & -3 & 4
\end{array}
\rightB^T \\
&=&
\leftB
\begin{array}{rrr}
\frac{1}{\sqrt{26}} & -\frac{3}{\sqrt{26}} & \frac{4}{\sqrt{26}}
\end{array}
\rightB^T
\end{eqnarray*}

Anda dapat memverifikasi dengan menggunakan Definisi \ref{def:lengthofvector} bahwa $\vectlength \vect{u} \vectlength = 1 $.
\end{solution}
